% Physical unit 019; accepted Session 02.
\GordanUnitStart{019}{printed page 14}
\GordanTopFolio{14}
\[
=(10)+{(11)\cos^2l\over \cos l_0\sqrt{1-e^2\cos^2l}}
\left\{{\alpha\sin l_0\cos\varphi+\beta{du\over d\varphi}
\over \alpha\sin l_0\sin\varphi+\beta u}
-{ {dv\over d\varphi}\over v}\right\}.
\]

\begin{center}\S\ 8.\end{center}

Quia in dextro latere aequationis quantitates:
\[
 (01)\quad (11)\quad \cos^2l\quad \tg l_0\quad {1\over \sqrt{1-e^2\cos^2l}}\quad
 \cos\varphi\quad {du\over d\varphi}\quad {dv\over d\varphi}\quad {1\over v}
\]
perpetuo finitae manent, functio $N$ hoc tantum modo infinita evadere potest, si:
\[
 \alpha\sin l_0\sin\varphi+\beta u=0
\]
fit.

Hinc sequitur, ut integrale
\[
 \int_{\psi_1}^{\psi_2}d\varphi\{(11)\delta l'\delta l'+2(01)\delta l\delta l'+(00)\delta l\delta l\}
\]
tamdiu positivum sit, quamdiu constantes $\alpha$ et $\beta$ ita determinari possint, ut
\[
 \alpha\sin l_0\sin\varphi+\beta u
\]
non evanescat, tamdiu igitur, quoad:
\[
 {u\over \sin l_0\sin\varphi}
\]
non omnes, inter $-\infty$ et $+\infty$ positos, numeros induat, nam si
\[
 {u\over \sin l_0\sin\varphi}
\]
aliquem valorem \gq{$a$} in hoc spatio non induit, posito:
\[
 \alpha+\beta a=0
\]
nunquam est:
\[
 \alpha\sin l_0\sin\varphi+\beta u=0.
\]

\begin{center}\S\ 9.\end{center}

Quod
\[
 {d\over d\varphi}\left\{ {u\over \sin l_0\sin\varphi}\right\}
 =-{1-\sin^2l_0\cos^2\varphi\over \sin l_0\cos l_0\sin^2\varphi}
\]
semper negativum est et posito $\psi=\psi_1$:
\[
 {u\over \sin l_0\sin\varphi}=0;
\]
$\psi$ crescente,
\[
 {u\over \sin l_0\sin\varphi}
\]
induit deinceps omnes valores a $0$ usque ad $-\infty$ et, relapsum a $-\infty$ ad $+\infty$, a $\infty$ ad $0$.
